% English translation of questions/5780/semA-moedC/q1.tex (metadata is taken from the Hebrew file).

\begin{question}
Investigate the function $y=\frac{2x-1}{(x-1)^2}$ according to the following items:
\begin{enumerate}
  \item The domain of the function, intersection points with the axes, evenness, oddness.
  \item Continuity, points of discontinuity.
  \item Asymptotes
  \item Intervals of increase and decrease, extremum points.
  \item Intervals of convexity and concavity, inflection points.
  \item Graphical description.
\end{enumerate}
\end{question}

\begin{hint}
When differentiating by the quotient rule, it is convenient to cancel one factor $(x-1)$ from the numerator and the denominator before simplifying.
\end{hint}

\begin{solution}
Denote $f(x)=\frac{2x-1}{(x-1)^2}$.
\begin{enumerate}
  \item \textbf{Domain:} $D=\{x\in\R: x\neq1\}$.

  \textbf{Intersection with the axes:} $f(0)=\frac{-1}{1}=-1$, i.e. $(0,-1)$ on the $y$-axis. $f(x)=0\iff 2x-1=0\iff x=\frac12$, i.e. $(\frac12,0)$ on the $x$-axis.

  \textbf{Evenness:} the domain is not symmetric with respect to $0$ ($-1\in D$ but $1\notin D$), hence the function is neither even nor odd.

  \item $f$ is a quotient of polynomials, hence continuous on its entire domain. At $x=1$: the numerator tends to $1>0$ and the denominator $(x-1)^2\to0^+$, hence $\lim_{x\to1^\pm}f(x)=+\infty$ --- a discontinuity of the second kind (infinite).

  \item \textbf{Vertical:} $x=1$ (from both sides $f\to+\infty$).
  \textbf{Horizontal:} $\lim_{x\to\pm\infty}\frac{2x-1}{(x-1)^2}=\lim_{x\to\pm\infty}\frac{\frac2x-\frac1{x^2}}{(1-\frac1x)^2}=0$, hence $y=0$ is a horizontal asymptote at $\pm\infty$ (and hence there is no oblique asymptote).

  \item By the quotient rule:
  \[ f'(x)=\frac{2(x-1)^2-(2x-1)\cdot2(x-1)}{(x-1)^4}=\frac{2(x-1)-2(2x-1)}{(x-1)^3}=\frac{-2x}{(x-1)^3}. \]
  $f'(x)=0\iff x=0$. Signs:
  \begin{itemize}
    \item $x<0$: numerator $-2x>0$, denominator $<0$, hence $f'<0$ --- decreasing.
    \item $0<x<1$: numerator $<0$, denominator $<0$, hence $f'>0$ --- increasing.
    \item $x>1$: numerator $<0$, denominator $>0$, hence $f'<0$ --- decreasing.
  \end{itemize}
  Hence $f$ is decreasing on $(-\infty,0)$ and on $(1,\infty)$ and increasing on $(0,1)$. At $x=0$ the derivative changes from negative to positive: a \textbf{local minimum} $(0,-1)$. (This is also an absolute minimum: for $x<\frac12$ the function decreases to $-1$ and then increases, and for $x\ge\frac12$, $f\ge0$.)

  \item Differentiate $f'(x)=-2x(x-1)^{-3}$:
  \[ f''(x)=-2(x-1)^{-3}+6x(x-1)^{-4}=\frac{-2(x-1)+6x}{(x-1)^4}=\frac{2(2x+1)}{(x-1)^4}. \]
  The denominator is positive for every $x\neq1$, hence the sign is that of $2x+1$:
  $f''<0$ for $x<-\frac12$ --- $f$ is concave ($\cap$) on $(-\infty,-\frac12)$;
  $f''>0$ for $-\frac12<x<1$ and $x>1$ --- $f$ is convex ($\cup$) on $(-\frac12,1)$ and on $(1,\infty)$.
  At $x=-\frac12$ the sign changes and the function is continuous, hence this is an \textbf{inflection point}: $f(-\frac12)=\frac{-2}{9/4}=-\frac89$, i.e. $(-\frac12,-\frac89)$.

  \item \textbf{Graphical description:} on the left the graph starts below the $x$-axis, close to the asymptote $y=0$ as $x\to-\infty$, decreases (concave) to the inflection point $(-\frac12,-\frac89)$, continues to decrease (convex) to the minimum $(0,-1)$, increases and crosses the $x$-axis at $(\frac12,0)$ and tends to $+\infty$ as $x\to1^-$. To the right of $x=1$ it decreases from $+\infty$ and approaches $y=0$ from above as $x\to\infty$ (convex).
\end{enumerate}
\end{solution}
